%! Simplifying Radicals; Adding \& Subtracting — Unit 6, Lesson 6.2 — Student Packet Cover Sheet
\documentclass[10pt]{article}
\usepackage{algebra2-article}
\usepackage{algebra2-boxes}
\usepackage{ltablex}
\keepXColumns

\begin{document}

% Full-bleed forest banner
\begin{tikzpicture}[remember picture, overlay]
  \fill[forest] (current page.north west)
    rectangle ([yshift=-0.9in]current page.north east);
\end{tikzpicture}
\vspace{-0.8in}

\begin{center}
  {\color{white}\textbf{\LARGE Algebra 2: Shepherd}}\\[4pt]
  {\color{forestmid}\large\textbf{Unit 6 \quad Radical Functions}}\\[6pt]
  {\color{white}\textbf{\large Lesson 6.2 \quad Simplifying Radicals; Adding \& Subtracting}}
\end{center}

\vspace{0.22in}
\namedateperiod
\vspace{0.18in}

\begin{learningtargetbox}
\small
\begin{enumerate}[leftmargin=1.5em, itemsep=5pt]
  \item \ldots{} use the \textbf{product} and \textbf{quotient} properties of radicals, and explain why they are just yesterday's exponent rules --- and why there is \textbf{no} such rule for addition.
  \item \ldots{} write a radical in \textbf{simplest form}: numeric radicands, algebraic radicands, and indices above $2$, by pulling out the \emph{largest} perfect $n$th power.
  \item \ldots{} \textbf{add and subtract} radicals by collecting \textbf{like radicals} --- simplifying first, because two radicals can be like radicals in disguise.
\end{enumerate}
\end{learningtargetbox}

\vspace{0.14in}

\begin{tocbox}
\small
\renewcommand{\arraystretch}{1.5}
\begin{tabularx}{\linewidth}{@{} c l X r @{}}
\rowcolor{forestbg}
\textbf{\#} & \textbf{Component} & \textbf{Description} & \textbf{Score} \\
\midrule
1 & Warm-Up        & Spiral review: perfect squares and largest factors, combining like terms, and the exponent rules & \blank{1.2cm} \\
2 & Guided Notes   & The product and quotient properties, simplest radical form, higher indices, algebraic radicands, and like radicals & \blank{1.2cm} \\
3 & Group Activity & \emph{Hiding in Plain Sight} --- simplifying, sorting into like-radical families, and a proof, by tier & \blank{1.2cm} \\
4 & Exit Ticket    & Quick check: simplify, combine, explain the addition error, plus an SOL-style item & \blank{1.2cm} \\
5 & Homework       & Mixed simplifying and combining, the quotient property, and a pendulum model & \blank{1.2cm} \\
\midrule
  & \multicolumn{2}{r}{\textbf{Total}} & \blank{1.2cm} \\
\end{tabularx}
\end{tocbox}

\vspace{0.10in}

\begin{remindbox}
\small A radical \textbf{splits over multiplication} and \textbf{never over addition}. The product
property $\sqrt[n]{ab}=\sqrt[n]{a}\cdot\sqrt[n]{b}$ is not a new rule --- it is yesterday's
$(ab)^{1/n}=a^{1/n}b^{1/n}$ in radical clothing. There is no matching rule for sums, and the
arithmetic proves it: $\sqrt{9}+\sqrt{16}=7$, while $\sqrt{9+16}=5$. To put a radical in
\textbf{simplest form}, pull out the \emph{largest} perfect $n$th power --- and remember that the
\textbf{index sets the group size}: pairs escape a square root, triples escape a cube root, so
$\sqrt[3]{9}$ is already finished even though $9$ is a perfect square. For variables, divide the
exponent by the index: the quotient goes outside, the remainder stays in
($\sqrt{x^{7}}=x^{3}\sqrt{x}$, because $\frac{7}{2}=3\frac{1}{2}$ --- simplifying a radical is
rewriting an improper fraction as a mixed number). Finally, radicals add exactly the way like terms
do: $3\sqrt5+7\sqrt5=10\sqrt5$ just as $3x+7x=10x$, and $3\sqrt5+7\sqrt2$ goes nowhere just as
$3x+7y$ does. The catch --- \textbf{simplify first, decide second}. $\sqrt{8}+\sqrt{18}$ looks
hopeless and is really $2\sqrt2+3\sqrt2=5\sqrt2$.
\end{remindbox}

\end{document}
